\documentclass[11pt,a4paper]{article}

% ==================================================
% VARIABLES – adjust per lab
% ==================================================
\newcommand{\course}{Software and Process Notations}
\newcommand{\labno}{1}
\newcommand{\labtopic}{Use Case Description}
\newcommand{\reportdate}{21 October 2026}
\newcommand{\headernames}{Stirbu, Hosseini, Stolzlechner}
\newcommand{\pdfauthorname}{Ricardo Stolzlechner}

% Laborantenliste: \laborant{Name}{E-Mail} (eine Zeile pro Person)
\newcommand{\members}{%
    \member{Alex Stirbu}{astirbu-blagut.cysb-m2026@fh-salzburg.ac.at}    
    \member{Manizheh Hosseini}{mhosseini.cysb-m2026@fh-salzburg.ac.at}
    \member{Ricardo Stolzlechner}{ricardo@stolzlechner.at}
}

% Verwendetes KI-Werkzeug (laut KI-Richtlinie 1.0, Kap. 1.7)
\newcommand{\aitool}{Claude Code (Anthropic, Claude Opus 5.5, token based billing)}

% --------------------------------------------------
% Packages
% --------------------------------------------------
\usepackage[utf8]{inputenc}
\usepackage[T1]{fontenc}
\usepackage[english]{babel}
\usepackage[a4paper,margin=2.5cm,headheight=14pt]{geometry}

\usepackage{amsmath, amssymb}
\usepackage{enumitem}
\usepackage{xcolor}
\usepackage{fancyhdr}
\usepackage{array}
\usepackage{booktabs}
\usepackage{graphicx}
\usepackage{listings}
\usepackage{setspace}
\usepackage{hyperref}

% Listings
\definecolor{codegreen}{rgb}{0,0.6,0}
\definecolor{codegray}{rgb}{0.5,0.5,0.5}
\definecolor{codepurple}{rgb}{0.58,0,0.82}
\definecolor{backcolour}{rgb}{0.95,0.95,0.92}

\lstdefinestyle{mystyle}{
    backgroundcolor=\color{backcolour},
    commentstyle=\color{codegreen},
    keywordstyle=\color{magenta},
    numberstyle=\tiny\color{codegray},
    stringstyle=\color{codepurple},
    basicstyle=\ttfamily\footnotesize,
    breaklines=true,
    captionpos=b,
    keepspaces=true,
    numbers=left,
    numbersep=5pt,
    showstringspaces=false,
    tabsize=2
}
\lstset{style=mystyle}

% Title page
\newcommand{\member}[2]{#1 & #2 \\}

% Spacing
\onehalfspacing
\setlength{\parindent}{0pt}
\setlength{\parskip}{6pt}

% --------------------------------------------------
% Header and footer
% --------------------------------------------------
\pagestyle{fancy}
\fancyhf{}
\lhead{\course}
\rhead{\headernames}
\cfoot{Page \thepage\ of \pageref{LastPage}}

% --------------------------------------------------
% Hyperref settings
% --------------------------------------------------
\hypersetup{
    colorlinks=true,
    linkcolor=blue,
    urlcolor=blue,
    pdftitle={\course\ -- Lab \labno: \labtopic},
    pdfauthor={\pdfauthorname}
}

% ==================================================
% Document
% ==================================================
\begin{document}

% --------------------------------------------------
% Page 1: Title page
% --------------------------------------------------
\pagenumbering{Alph}
\begin{titlepage}
    \centering
    \vspace*{4cm}
    {\LARGE \textbf{Lab Report - Lab Unit \labno}}\\[0.5em]
    {\Large \labtopic}\\[0.5em]
    {\large \course}

    \vspace{3cm}

    \begin{tabular}{l l}
        \toprule
        \textbf{Lab Member} & \textbf{E-mail} \\
        \midrule
        \members
        \bottomrule
    \end{tabular}

    \vfill
    {\large \reportdate}
\end{titlepage}
\pagenumbering{arabic}

% --------------------------------------------------
% Page 2: Content
% --------------------------------------------------
\subsection*{Use Case Name}

Automated Valet Parking \& Charging (Compliance and Data Privacy
Management)

\subsection*{Description}

Customers arrive at the entry terminal, the license plate is recognized and
a session is created. The driver now has the choice to bind their identity
to the session via a bank card through a payment token. If the driver
cannot or does not want to use a card, a printed ticket with a QR code
and an alphanumeric code is issued instead. The driver can optionally
request charging, and the car is then parked automatically.

\subsection*{Preconditions}

\begin{itemize}
    \item The facility is operational with at least one free space (and a free
charger if charging is requested).
\item The entry and exit terminals are functional, with a working card reader,
a connection to the payment provider and ticket paper available.
\item The vehicle is a passenger car within the size and weight limits.
\item The privacy notice and support hours are visible at the terminal.
\end{itemize}


\subsection*{Postconditions}

\begin{itemize}
    \item The car was parked, charged if requested, and handed over to the
identified driver.
\item Payment is settled.
\item The plate, card token or ticket token hash, and all other session
identifiers are deleted. Only an anonymous billing record (session ID,
timestamps, amount, kWh) remains for the legally required period.
\item Anonymized data is available for statistics.
\end{itemize}

\subsection*{Frequency of Use}

Continuous throughout the day, with peaks at commuting hours and events.

\subsection*{Normal Course of Events}

\begin{enumerate}
    \item The driver drives into the entry bay and stops.
    \item The system checks that the vehicle is eligible (type and size).
    \item The system captures the plate image, extracts the plate text and
    immediately discards the image.
    \item The driver confirms the recognized plate.
    \item The system displays the tariff, support hours and a short privacy
    notice.
    \item The driver optionally selects charging.
    \item The driver taps or inserts a bank card.
    \item The system requests a payment token (and a pre-authorization
    reference if applicable) from the payment provider and stores only
    the token.
    \item The system creates the session, linking plate and card token.
    \item The driver leaves the bay, and the system verifies that it is empty.
    \item The system assigns a space, parks the car and starts charging if
    requested. The session status becomes ``parked''.
    \item The driver taps the same card at the exit terminal.
    \item The system matches the card token to the active session.
    \item The system brings the car to the handover bay and stops charging.
    \item The system recognizes the plate again and matches it with the
    session as a second check.
    \item The system calculates the fee and settles payment with the card.
    \item The driver takes the car.
    \item The system closes the session and deletes plate and card token,
    keeping only the anonymous billing record.
\end{enumerate}

\subsection*{Alternate Course}

\paragraph{Ticket instead of card (fallback)}
In step 7 the driver chooses ``no card'' (for example cash customer,
foreign cards, or a preference for anonymity). The system generates a
random token, stores only its hash, creates the session, and prints a
ticket (QR code, alphanumeric code, ticket number, entry time,
support hours). The course then continues at step 12.

\paragraph{Retrieval with ticket}
In step 12 the driver scans the QR code. If it is unreadable, the driver
types the alphanumeric code together with the plate. The course
continues at step 14.

\paragraph{Driver leaves immediately after check-in}
The system cancels the session, releases any pre-authorization, and
deletes all data. Nothing is charged (or a short free grace period
applies).

\paragraph{Driver rejects the recognized plate (step 4)}
The driver corrects it manually or aborts. If aborted, all data is
deleted.

\paragraph{Charging not requested}
The charging steps are skipped.

\subsection*{Exceptions}

\paragraph{Plate not recognized}
The system retries up to 3 times, then offers manual plate entry with
confirmation. The system compares the typed plate with the
captured image and calculates a confidence value. If the confidence
is at or above the threshold, the plate is accepted and the course
continues at step 4 of the normal course. If the confidence is below
the threshold, the system refuses entry, discards the image, and
instructs the driver to leave the entry bay. The driver can clean the
plate or fix the obstruction outside and try again. No session is
created and nothing is stored.

\paragraph{Vehicle not eligible (for example a motorcycle or oversized vehicle)}
Entry is refused before any session is created, and nothing is stored.

\paragraph{No free space or charger}
The driver is informed before a session is created. Parking without
charging may be offered.

\paragraph{Card declined or not readable at entry}
The system allows one more attempt, then offers another card or the
ticket path.

\paragraph{Payment provider unavailable}
The system offers the ticket path so entry remains possible. Payment
is settled at exit. If the ticket path is also unavailable, entry is refused
and the driver leaves the bay.

\paragraph{Ticket printer out of paper or defective (ticket path only)}
No session is created, and the driver is offered the card path. The
system sends a maintenance alert to the operator.

\paragraph{Obstruction or person detected in the bay}
Parking is paused until the bay is clear.

\paragraph{Plate mismatch at handover (step 15)}
The car is not released, and the driver is directed to support.

\paragraph{Too many failed attempts at the exit terminal}
The session is temporarily locked and the driver is directed to
support.

\paragraph{Card and ticket lost, replaced or unusable}
The terminal tells the driver to contact support and shows the
support hours. At the desk, staff visually check the original ID and the
original registration document, and confirm that the plate matches
the parked car. If the registered holder differs from the driver, an
authorization letter is required or the release is declined. Staff then
release the session manually, and the plate match at the handover
bay still applies. Nothing is scanned or copied, and the log records
only session reference, time, staff member, documents checked and
outcome.

\paragraph{System or terminal failure}
No session is created, and existing sessions stay intact. At entry, the
system refuses entry and instructs the driver to leave the bay. At the
exit terminal, the driver is directed to support.

\subsection*{Includes}

\begin{itemize}
    \item Recognize License Plate (entry and handover)
    \item Bind Payment Card to Session
    \item Park Vehicle Automatically in Facility
    \item Charge Vehicle
    \item Settle Payment
    \item Retrieve Vehicle
\end{itemize}

\subsection*{Special Requirements}

\paragraph{GDPR}

\begin{itemize}
    \item The plate image is held only for the duration of plate extraction
    and, if needed, the manual-entry confidence check. It is
    discarded immediately afterwards, whether the plate is accepted
    or rejected.
    \item The plate is processed only for parking, charging and billing
    (lawful basis: performance of the contract). The identity check
    by customer support when card and ticket are
    lost rests on legitimate interest, backed by a documented balancing
    test.
    \item Signage and a privacy notice appear before capture.
    \item Images are discarded right after plate extraction.
    \item Only the payment token is stored, never the card number. The
    payment provider is a processor and needs a data processing
    agreement.
    \item Ticket tokens are random, never derived from the plate, and the
    QR code contains no personal data.
    \item Plate and token are deleted at session closure, and billing records
    contain no plate data.
    \item Data is stored on the company server with access control and
    encryption, and backups follow the deletion rules.
    \item No biometric processing, and no copies of identity documents.
    \item Statistics use only anonymized or aggregated data.
    \item A DPIA (Data Protection Impact Assessment (Art. 35 GDPR)) is
    carried out before operation. That is a risk analysis regarding
    systematic monitoring of people recorded on the facility,
    automated recognition technology and large-scale processing.
\end{itemize}

\paragraph{Other}

\begin{itemize}
    \item No registration is required, and both card and ticket paths are
    always offered.
    \item Card handling follows PCI DSS (Payment Card Industry Data
    Security Standard), by delegating card data entirely to the
    payment provider.
    \item Retrieval with a valid card or ticket is available 24/7. Lost
    credential recovery is available only during support hours, which
    are shown on the terminal and tickets.
    \item The alphanumeric code avoids easily confused characters (0/O,
    1/I/l).
    \item Session creation completes within a defined time (for example
    under 10 seconds).
\end{itemize}

\subsection*{Assumptions}

\begin{itemize}
    \item The entry area allows a driver to turn around or leave without a
    session, so a refused driver is never stuck in the bay.
    \item The driver has a passenger car with a valid, legible license plate
    (motorcycles and other vehicle types are out of scope).
    \item The car fits the size and weight limits of the facility.
    \item The driver has a bank card or accepts the ticket path, and keeps the
    card or ticket for the whole stay.
    \item For recovery, the driver can come to the desk in person with ID and
    registration document. If the papers are locked in the parked car,
    staff may fetch them under supervision after the ID check.
    \item The driver accepts the tariff and terms at the terminal.
\end{itemize}

% --------------------------------------------------
% Page 3: AI disclosure
% --------------------------------------------------
\clearpage
\section*{Statement on the Use of AI}
\addcontentsline{toc}{section}{Statement on the Use of AI}

In preparing this lab report, \aitool{} was used as a generative AI tool
exclusively to assist with the following tasks:

\begin{itemize}
    \item creating and designing the \LaTeX{} layout,
    \item checking spelling and grammar,
    \item research.
\end{itemize}

The \LaTeX{} template of this report, including this statement, was created
with the assistance of the AI tool and reviewed and adapted by the lab members.
Apart from that, all content, results and conclusions were developed
independently; no AI-generated text was adopted. The tool was used in line with
the principles of \emph{responsible AI use} (transparency, verifiability, human
oversight). Sources found through research were checked for accuracy
independently. All content was reviewed, validated and finally approved by the
lab members (cf.\ FH Salzburg AI guideline \emph{KI-Richtlinie}, version~1.0,
ch.~1.3 and 2.1.2).

\label{LastPage}
\end{document}
